\documentclass[tikz,border=8pt]{standalone}
% 公共导言：XeLaTeX + TikZ + 中文字体
\usepackage{xeCJK}
% 优先使用开源的 Noto Sans SC（OFL）；没有时回退到系统字体，便于他机复现。
\IfFontExistsTF{Noto Sans SC}
  {\setCJKmainfont{Noto Sans SC}}
  {\IfFontExistsTF{Source Han Sans SC}
    {\setCJKmainfont{Source Han Sans SC}}
    {\setCJKmainfont{Microsoft YaHei}}}
\usepackage{amsmath}
\usetikzlibrary{positioning,arrows.meta,shapes.geometric,fit,backgrounds,calc,decorations.pathreplacing}
\tikzset{
  block/.style={draw,rounded corners=2pt,align=center,font=\small,inner sep=4pt},
  existing/.style={block,fill=blue!8,draw=blue!55!black},
  planned/.style={block,fill=gray!10,draw=gray!60!black,dashed},
  ui/.style={block,fill=green!8,draw=green!45!black},
  ext/.style={block,fill=orange!10,draw=orange!65!black},
  auth/.style={block,fill=red!7,draw=red!55!black},
  layer/.style={draw,rounded corners=4pt,fill=black!3,inner sep=6pt},
  arr/.style={-{Latex[length=2mm]},thick},
  darr/.style={-{Latex[length=2mm]},thick,dashed},
  lbl/.style={font=\scriptsize,fill=white,inner sep=1pt}
}

\begin{document}
\begin{tikzpicture}[x=1cm,y=1cm,
  dev/.style={core, minimum width=4.60cm, minimum height=5.40cm},
  peer/.style={ext, minimum width=3.90cm, minimum height=2.10cm},
  spec/.style={ui, minimum width=3.90cm, minimum height=2.10cm},
  up/.style={-{Latex[length=2.2mm]}, draw=cblue, line width=1.0pt},
  down/.style={-{Latex[length=2.2mm]}, draw=corange, line width=1.0pt}]

  % 两块客板放在左侧，房主独占右侧
  \node[peer] (opp) at (-4.40,2.40)
    {\textbf{玩家板}\\[1pt]
     \scriptsize 通过蓝牙连到房主\\[-1pt]
     \scriptsize 只发送自己的操作\\[-1pt]
     \scriptsize 等房主发回棋盘后更新\\[-1pt]
     \textcolor{corange}{\scriptsize 不会先改自己的棋盘}};
  \node[spec] (spec) at (-4.40,-2.00)
    {\textbf{旁观板}\\[1pt]
     \scriptsize 通过蓝牙连到房主\\[-1pt]
     \scriptsize 只能看，不能下棋\\[-1pt]
     \scriptsize 能看到双方棋盘和预览};
  \node[dev] (host) at (3.60,0.20)
    {\textbf{房主板}\\[2pt]
     \scriptsize 同时是玩家和裁判\\[2pt]
     \scriptsize 房主也是蓝牙服务端\\[6pt]
     \scriptsize 保存唯一的棋局\\[-1pt]
     \scriptsize 判断每一步是否合法\\[-1pt]
     \scriptsize 结算后把棋盘发给各板\\[-1pt]
     \scriptsize 把预览转给旁观者\\[4pt]
     \textcolor{cblue}{\scriptsize 棋局以房主为准}};

  % 玩家板 <-> 房主
  \draw[up]   (opp.east |- {0,2.65}) -- node[lab, above]{加入房间、发送操作} (host.west |- {0,2.65});
  \draw[down] (host.west |- {0,2.15}) -- node[lab, below]{房间状态、棋盘、结果} (opp.east |- {0,2.15});

  % 旁观板 <-> 房主
  \draw[up]   (spec.east |- {0,-1.75}) -- node[lab, above]{加入旁观} (host.west |- {0,-1.75});
  \draw[down] (host.west |- {0,-2.25}) -- node[lab, below]{房间状态、棋盘、预览} (spec.east |- {0,-2.25});

  % 说明
  \node[draw=ink!20, fill=ink!3, rounded corners=3pt, align=left,
        font=\scriptsize, text=ink!85, inner sep=5pt, text width=12.4cm]
    at (0.00,-4.65)
    {找房间：房主会不断广播自己的房间，其他板在列表里看到后加入，不需要中心服务器。\\
     预览：只有自己看得到；房主会转给旁观者，对手看不到。\\
     掉线后：10 秒内可以回到原来的位置；等待阶段房主退出时，另一位玩家接管房间，旁观者跟着新房主。};

\end{tikzpicture}
\end{document}
